% !TeX program = xelatex
% business-justification.tex
% Document ENG-RPT-02 - Business Justification: On-Premises ML Platform
%
% Build (from this folder, requires XeLaTeX; shared/intake-style.sty is a copy of the
% engineering_doc_templates version, kept in this repo so it builds on its own):
%   xelatex -interaction=nonstopmode business-justification.tex
%   xelatex -interaction=nonstopmode business-justification.tex   (rerun for TOC/page refs)
%
% Static report (no fillable fields). Shares branding, page layout, cover page
% and header/footer markings with the engineering intake template family via
% intake-style.sty. Companion document: ENG-RPT-01 Analysis of Alternatives.
%
% Main file: preamble, cover page, document control and the section list. Each
% section is its own file in sections/ (edit those). Table column weights (Y{...})
% in each table must sum to the number of Y columns so the table fills \textwidth.

\documentclass[11pt]{article}

\usepackage{../../shared/intake-style}
\usepackage{../../shared/report-components}

\SetDocIdentity{Business Justification: On-Premises ML Platform}{ENG-RPT-02}{v0.4-draft}

% Replace these placeholders once branding and markings are confirmed.
\SetCompanyName{[Engineering Center Name]}
\SetMarkingText{FOR OFFICIAL USE ONLY}

\SetDocPurpose{%
This report makes the business case for building a two-computer, on-premises
system that can be used with Controlled Unclassified Information (CUI):
Computer~1, a network-attached storage (NAS) and private-cloud server for file
sharing, storage, software licensing and the application tier; and Computer~2,
one or two machine-learning (ML) nodes that serve a large-context large language
model (LLM) to many users at once and run fine-tuning, with the second node
providing training capacity and backup. It compares that investment with
continuing to rely on the two AI tools currently provided, and with a compliant
cloud service. Dollar figures that depend on measured usage are illustrative and
must be replaced before the report is used to request funding. The companion
Analysis of Alternatives (ENG-RPT-01) gives the technical design and sizing.%
}

\begin{document}

\CoverPage

\TodoAnchor{bj-ctrl}
\ReportControlBlock{Draft}{v0.4}{2026-10-22}{\PH{Author}}{ENG-RPT-01}{\PH{Recipients}}

\RevisionHistory{%
  v0.1 & \PH{Author} & 2026-10-14 & Initial business justification (large-budget single ML platform). \\
  v0.2 & \PH{Author} & 2026-10-16 & Re-based to \$15K--\$55K budget; \$100K salary and 2~hr/week upkeep. \\
  v0.3 & \PH{Author} & 2026-10-22 & Two-computer design (NAS + ML node[s]); dual-node training/failover value; Computer~1 benefits; LaTeX format. \\
  v0.4 & \PH{Author} & 2026-10-22 & Sponsor answers: $\geq$20 named / 12 simultaneous users, archive retrieval and training, Dell-configured pricing; scenarios recomputed. \\
}

\begin{ReportNote}{Estimates, not quotes}
Figures marked \emph{illustrative} are placeholders that show the calculation
method. Replace them with measured values (users, hours, labor rate, license
cost) before this report is used to request funding. Items in \PH{brackets} still
need a value. This document contains no CUI; verify that before adding
project-specific detail.
\end{ReportNote}

% ===========================================================================
\IntakeSection{Items To Complete Before Release}

\Guidance{Every placeholder and open decision in this report is listed here and links to the
place it appears. Resolve each item, tick the Status box, and update the revision history
before the report is sent outside the team. ``Required'' items are missing information;
``Confirm'' items are assumptions that need a named person to agree to them.}\\[4pt]

{\footnotesize
\begin{xltabular}{\textwidth}{@{}c Y{1.5} Y{0.7} c c Y{0.8}@{}}
  \toprule
  \# & Item & Needed from & Priority & Status & Go to \\
  \midrule
  \endhead
  1 & Document identity and handling: author, linked intake reference, recipients, center name and logo, marking text, version and date. & Document owner, security office & Required & $\square$ & \TodoLink{bj-ctrl}{Control block} \\
  2 & Organization name and the names of Tool~A and Tool~B. & Sponsor & Required & $\square$ & \TodoLink{bj-purpose}{Purpose} \\
  3 & Existing-tools table: hosting, CUI authorization, model, context window, output and upload limits, API, retrieval, rate limits, cost model, retention; plus usage statistics. & IT, tool owners & Required & $\square$ & \TodoLink{bj-tools}{Existing tools} \\
  4 & Current arrangement and cost for file storage, license servers and private-cloud software. & IT & Required & $\square$ & \TodoLink{bj-arrangement}{Pain points} \\
  5 & Status-quo costs: current license, subscription and support fees, contract end dates, and any vendor upgrade quotes. & Finance, contracts & Required & $\square$ & \TodoLink{bj-statuscost}{Costs} \\
  6 & Dell quotes for Computer~1 and the nodes (replace the 1.3 multiplier); Windows Server, database and client-access licensing (about \$2--8K, not yet included). & Dell account team, IT & Required & $\square$ & \TodoLink{bj-costs}{Costs} \\
  7 & Labor assumptions (hourly rate, upkeep hours, setup hours) still accepted. & Sponsor, finance & Confirm & $\square$ & \TodoLink{bj-labor}{Labor assumptions} \\
  8 & Replace the illustrative scenario inputs with measured values: active users by team, roles and rates, hours saved (time study and pilot), realization factor. & Sponsor, pilot team & Required & $\square$ & \TodoLink{bj-scen}{Scenarios} \\
  9 & Decide D3, D4 or S2 once the training workload and the required recovery time are known; update the second-node analysis. & Sponsor & Required & $\square$ & \TodoLink{bj-nodeb}{Second ML node} \\
  10 & Workflow list: users per workflow, domain for the fine-tuned assistant, current arrangement for hosted software. & Department leads & Required & $\square$ & \TodoLink{bj-usecases}{Use cases} \\
  11 & Baseline data: usage statistics, time study, unapproved-tool survey, archive size, fine-tuning workload, operator time. & Project team & Required & $\square$ & \TodoLink{bj-baseline}{Baseline data} \\
  12 & Answers to the open questions: archive size and formats, standards licenses, Autodesk and Bluebeam licensing, server standard, failover target, backup copy, cloud fallback, branding. & Sponsor & Required & $\square$ & \TodoLink{bj-open}{Open questions} \\
  13 & Approvers named, decisions and dates recorded. & Sponsor & Required & $\square$ & \TodoLink{bj-approve}{Approvals} \\
  14 & Cross-check all figures against ENG-RPT-01 and update both documents together. & Document owner & Required & $\square$ & \TodoLink{bj-costs}{Costs} \\
  15 & No bracketed placeholders remain in the finished PDF; compiled twice so the contents and cross-references are current; revision history updated. & Document owner & Required & $\square$ & --- \\
  \bottomrule
\end{xltabular}}

% ===========================================================================
\ReportSection[sec:labor]{Labor Assumptions Used Throughout}
\TodoAnchor{bj-labor}

\TableTitle{Labor assumptions (change here if they differ)}
{\small
\begin{tabularx}{\textwidth}{@{}Y{0.857}Y{0.771}Y{1.372}@{}}
  \toprule
  Assumption & Value & Basis \\
  \midrule
  Operator salary & \$100,000/yr & Given. \\
  Hourly rate & \textbf{\$48.08/hr} & \$100,000 $\div$ 2,080~hr; salary only, no fringe or overhead multiplier. \\
  Upkeep, ML node(s) & \textbf{2~hr/week} for one node, 3~hr/week with two & 2~hr/week given; 1~hr/week for a second node is assumed. \\
  Upkeep, Computer~1 & 1~hr/week & Assumed; accepted by the sponsor for now. \\
  One-time setup labor & Computer~1: 80--160~hr; first ML node: 160--320~hr; second ML node: 40--80~hr & Assumed; accepted by the sponsor for now. \\
  Rate used for users' saved time & \$48.08/hr & Same salary-only basis, for consistency. \\
  Users & At least 20 named; about 12 at the same time & Given. \\
  Hardware pricing basis & Street price $\times$ 1.3 (range 1.15--1.4) & Approximates Dell-configured servers until a quote is received (Dell is the stated vendor). \\
  \bottomrule
\end{tabularx}}

% ===========================================================================
\ReportSection[sec:purpose]{Purpose}
\TodoAnchor{bj-purpose}

Justify an investment in a dedicated, on-premises system authorized for CUI,
consisting of a storage and private-cloud computer plus one or two ML nodes,
compared with continued reliance on the two AI tools currently provided by
\PH{Organization} (``Tool A'' and ``Tool B'').

% ===========================================================================
\ReportSection[sec:exec]{Executive Summary}

\begin{itemize}
  \item The current tools are \textbf{functionally limited and constrained by small
    context windows}, so they cannot support the document-heavy, multi-step
    engineering workflows the organization needs (long specifications,
    multi-document comparison, large scripts and codebases, report generation).
  \item Staff either work around the limits by hand, abandon the tools, or are
    tempted to use unapproved external tools, which is a CUI spill risk.
  \item A self-hosted ML node provides \textbf{128K-token context}, an \textbf{API for
    automation}, \textbf{retrieval over internal repositories}, \textbf{multi-user
    concurrency} and \textbf{full control of CUI} inside the authorization boundary.
  \item A \textbf{second ML node} lets fine-tuning run \emph{at the same time} as
    production use without slowing users down, gives a failover path, and lets
    upgrades be tested before they reach users.
  \item \textbf{Computer~1} (about \$15K at street prices, about \$20K Dell-configured)
    consolidates file storage, backups,
    software license servers and a private cloud for specific software. It is needed
    whichever LLM option is chosen, so the cost that depends on the LLM decision is
    the ML node(s).
  \item \textbf{Recommended configuration D3:} Computer~1 + a 2-GPU serving node + a
    1-GPU training node that can also serve as a second node. Hardware is about \$64K at
    street prices and about \$83K Dell-configured (range \$66--101K). Three-year cost
    including labor is about \$131--192K Dell-configured, of which the ML side is about
    \$98--145K (about \$123K expected). A 1-GPU serving node is rated for 3--8
    simultaneous users, below the stated 12, so the single-GPU configurations S1 and D1
    are not recommended.
  \item \textbf{The NAS archive becomes usable knowledge.} The model can answer questions
    from archived project data with citations (retrieval) and can be fine-tuned on
    selected documents, for example a licensed copy of NFPA~780. Retrieval is the primary
    mechanism; fine-tuning teaches vocabulary, formats and tasks.
  \item \emph{Illustrative} result on the ML side alone, for 30 active users saving
    2~hr/week: payback in about \textbf{15 months} and a 3-year return of about 94\%.
    Attributing Computer~1 to the project lengthens payback to about 21 months. If only
    the minimum 20 users each save 1~hr/week, time savings alone do \emph{not} repay the
    cost. The 3-year cost is recovered at about \textbf{1~hour per user per week}
    (Section~\ref{sec:fin}).
  \item The strongest justification is not only cost: it is \textbf{workflows that
    are impossible today}, \textbf{lower data-handling risk}, and the ability to
    \textbf{adapt models to internal data} on equipment the organization controls.
\end{itemize}

% ===========================================================================
\ReportSection[sec:current]{Current State}

\ReportSubsection{Existing tools (complete the table)}
\TodoAnchor{bj-tools}

\TableTitle{Existing tools versus what the target workflows need}
{\footnotesize
\begin{xltabular}{\textwidth}{@{}Y{1.053}Y{0.842}Y{0.842}Y{1.263}@{}}
  \toprule
  Attribute & Tool A: \PH{name} & Tool B: \PH{name} & Needed for target workflows \\
  \midrule
  \endhead
  Hosting / owner & \PH{ } & \PH{ } & On-premises or accredited enclave \\
  CUI authorized? & \PH{ } & \PH{ } & Yes (required) \\
  Underlying model & \PH{ } & \PH{ } & Selectable, upgradable, auditable \\
  Max context window & \PH{e.g., 4K--32K} & \PH{ } & 128K tokens or more \\
  Max output length & \PH{ } & \PH{ } & 4K--8K tokens or more \\
  File and document upload limits & \PH{ } & \PH{ } & Large PDFs, spreadsheets, code \\
  API / automation & \PH{ } & \PH{ } & Yes (OpenAI-compatible) \\
  Retrieval over internal data & \PH{ } & \PH{ } & Yes, with per-group access control \\
  Concurrency / rate limits & \PH{ } & \PH{ } & Supports team-wide peak use \\
  Customization / fine-tuning & \PH{ } & \PH{ } & Prompt libraries, custom tools, adapter fine-tuning \\
  Code and script assistance quality & \PH{ } & \PH{ } & Strong \\
  Audit logging & \PH{ } & \PH{ } & Per-user, exportable to the SIEM \\
  Cost model & \PH{ } & \PH{ } & Predictable \\
  Roadmap control & Vendor & Vendor & Organization-controlled \\
  Retention and training on inputs & \PH{ } & \PH{ } & None; organization-controlled \\
  \bottomrule
\end{xltabular}}

\ReportSubsection{Observed pain points (validate with a user survey)}
\TodoAnchor{bj-arrangement}
\begin{enumerate}
  \item \textbf{Context limits} force document splitting, losing cross-references and
    giving inconsistent answers.
  \item \textbf{Limited functionality:} weak file handling, no custom tools, no API or
    workflow chaining.
  \item \textbf{Quality gaps} on technical reasoning, code and long structured output.
  \item \textbf{Low adoption and trust:} users stop using tools after poor results, so
    existing spend yields little value.
  \item \textbf{Shadow-AI risk:} users may paste work content into unapproved tools.
  \item \textbf{No adaptability} to organizational terminology, templates and standards.
  \item \textbf{No shared storage, license or private-cloud capability} tied to the
    AI tools: engineering data and software services are managed separately
    \PH{describe current arrangement}.
\end{enumerate}

% ===========================================================================
\ReportSection[sec:solution]{Proposed Solution (Summary)}

\textbf{Computer~1 -- NAS and private-cloud server.} Departmental file share,
bulk storage and backups, software license servers, a private cloud for specific
software, and the LLM application tier (gateway, chat interface, retrieval, vector
database, model registry, logging).

\textbf{Computer~2 -- ML node(s).} Open-weight LLMs (for example a roughly 120B-class
mixture-of-experts model) on workstation-class 96~GB GPUs, with an OpenAI-compatible
API, retrieval over controlled repositories, single sign-on and audit logging,
all inside the CUI boundary. Either one node, or two nodes: Node~A serves users while
Node~B trains and experiments, and can take over serving if Node~A fails. All ML
nodes use the same chassis class, so adding a third node later is a copy of an
existing one.

\TableTitle{Configurations (hardware only; see ENG-RPT-01 for detail)}
{\footnotesize
\begin{xltabular}{\textwidth}{@{}l Y{1.219} r Y{0.75} Y{1.031}@{}}
  \toprule
  Config. & Composition & Hardware & Active users* & Training and resilience \\
  \midrule
  \endhead
  S1 & Computer~1 + 1-GPU node & \$35K & 3--8 & Training shares the GPU; no failover. \\
  S2 & Computer~1 + 2-GPU node & \$44K & 8--20 & Split GPUs between serving and training; no failover. \\
  D1 & Computer~1 + two 1-GPU nodes & \$55K & 3--8 & Dedicated training node; full-capacity failover. \\
  \textbf{D3} & \textbf{Computer~1 + 2-GPU node (A) + 1-GPU node (B)} & \textbf{\$64K} & \textbf{8--20} & \textbf{Dedicated training node; reduced-capacity failover.} \\
  S4 & Computer~1 + 4-GPU node & \$65K & 25--50 & Single node; no failover. \\
  D4 & Computer~1 + two 2-GPU nodes & \$73K & 8--20 (A), more with both serving & Dedicated training node; full-capacity failover. \\
  \bottomrule
\end{xltabular}}
{\footnotesize * Simultaneous active users on the serving node(s). The stated need is at least 20 named users with about 12 at once, so S1 and D1 (1-GPU serving node) fall short. Hardware is at street prices; Dell-configured prices are about 1.15--1.4 times higher. Estimates to be confirmed by benchmark.}

\ReportSubsection{What to expect from the system}

\TableTitle{Expectations in plain terms (planning targets; ENG-RPT-01 gives the detail)}
{\footnotesize
\begin{xltabular}{\textwidth}{@{}Y{0.7} Y{1.3}@{}}
  \toprule
  Topic & What to expect \\
  \midrule
  \endhead
  Who can use it & At least 20 named users, with about 12 able to get answers at the same moment. Others wait briefly at peak. \\
  How much it can read at once & About 300 pages (128K tokens) in one request. Larger archives are searched by retrieval rather than read whole. \\
  Speed & First words in under 3~seconds for ordinary questions, then faster than reading speed. Very long documents take longer. \\
  Answers from the archive & Cited to the document, clause and page. An answer with no citation should not be relied on. \\
  Accuracy & Good first drafts and summaries; it can be wrong. Engineers review anything that goes into a deliverable. \\
  Learning your material & Retrieval makes archived documents and standards (for example NFPA~780) usable immediately; fine-tuning adapts style and vocabulary but does not replace retrieval. \\
  Where data goes & Nothing leaves the organization's boundary; everything is logged per user. \\
  If something fails & With two ML nodes, serving moves to Node~B at reduced capacity in D3; with one node, service stops until repair. Computer~1 relies on disk redundancy only. \\
  What it does not do & Read drawings as images, browse the internet, or train a new model from scratch. \\
  \bottomrule
\end{xltabular}}

% ===========================================================================
\ReportSection[sec:benefits]{Benefits}

\ReportSubsection{Capability benefits (what the existing tools cannot do)}

\TableTitle{Capabilities and example workflows}
{\footnotesize
\begin{xltabular}{\textwidth}{@{}Y{0.844}Y{0.938}Y{1.218}@{}}
  \toprule
  Capability & Business value & Example workflow \\
  \midrule
  \endhead
  \textbf{Large context (128K tokens, about 300 pages)} & Analyze whole specifications and reports in one pass. & Compare a 200-page specification to a test report and list gaps. \\
  \textbf{Retrieval with citations} & Find and cite answers across thousands of documents. & ``What did past programs decide about connector X?'' with sources. \\
  \textbf{API access} & Embed the LLM in scripts, pipelines and spreadsheets. & Nightly summary of new anomaly reports. \\
  \textbf{Code and script assistance on larger codebases} & Faster development, documentation and test generation. & Review a multi-file Python/MATLAB toolchain and generate unit tests. \\
  \textbf{Concurrent multi-user access with quotas} & Team-wide availability and predictable response. & 15 engineers using it at once at peak. \\
  \textbf{Batch processing} & Bulk tasks run overnight. & Classify and summarize 10,000 legacy records. \\
  \textbf{Expert assistants built from the NAS archive} & Retrieval over archived projects and licensed standards with clause-level citations; optional fine-tuning on reviewed examples. & Ask questions of NFPA~780 and past project reports and get cited answers. \\
  \textbf{Fine-tuning on internal data (Node~B)} & Learn terminology, templates and formats without slowing users. & Adapter tuned on past reports to draft in the standard format. \\
  \textbf{Model choice and upgrades} & Swap in better models as released; no vendor lock-in. & Adopt a new approved model after testing on Node~B. \\
  \textbf{Structured output and tool use} & Reliable JSON and forms. & Auto-populate traceability matrices. \\
  \bottomrule
\end{xltabular}}

\ReportSubsection{Benefits of the second ML node}
\begin{itemize}
  \item \textbf{Training without disruption.} Fine-tuning and evaluation jobs can run for
    hours or days on Node~B while Node~A keeps serving at full speed. On a single node,
    training must be scheduled off-hours or must take GPUs away from users.
  \item \textbf{Failover.} If Node~A fails, the gateway routes to Node~B (at reduced
    capacity in D3, full capacity in D1 and D4), so a hardware failure does not stop all
    LLM use.
  \item \textbf{Second serving node.} When no training job is running, Node~B can load the
    same model and share the load. Training takes priority, so this is spare capacity, not
    guaranteed capacity.
  \item \textbf{Safe upgrades.} Drivers, inference software and new models are tested on
    Node~B before production sees them, and patching can be done one node at a time.
  \item \textbf{Growth.} Because both nodes are the same class of machine, either can be
    given more GPUs (about \$9K each) and a third node can be added by repeating the build.
\end{itemize}
Section~\ref{sec:node-b} estimates whether these benefits justify the cost of Node~B.

\ReportSubsection{Benefits of Computer 1}
\begin{itemize}
  \item \textbf{One governed home for engineering data:} file share with access
    control, snapshots, quotas and tested backups, inside the CUI boundary.
  \item \textbf{License servers and a private cloud} for Autodesk products, Bluebeam and
    \PH{other software}, replacing ad hoc workstations or third-party hosting
    \PH{describe current arrangement and cost}. Whether a local license server is needed
    depends on the licensing model (ENG-RPT-01, Section on hosting Autodesk and Bluebeam).
  \item \textbf{Platform for the LLM application tier:} retrieval indexes, model and
    artifact registry and logs sit on protected, backed-up storage.
  \item \textbf{Independent value:} Computer~1 would be justified even if no LLM were
    deployed; the LLM adds a modest increment to its cost.
\end{itemize}

\ReportSubsection{Productivity benefits}
Published studies show wide ranges; validate in the pilot.

\TableTitle{Expected time reduction by task category}
{\small
\begin{tabularx}{\textwidth}{@{}Y{1.448}Y{0.517}Y{1.035}@{}}
  \toprule
  Task category & Baseline effort & Expected reduction \\
  \midrule
  Document summarization and review & 1--4~hr & 40--70\% (largest with long context) \\
  First-draft reports, emails, procedures & 1--3~hr & 30--50\% \\
  Information search across repositories & 0.5--2~hr & 50--80\% (needs retrieval) \\
  Code and script writing, debugging, tests & varies & 20--40\% \\
  Data extraction and transformation & 1--6~hr & 40--80\% \\
  Requirements and traceability analysis & 4--16~hr & 30--60\% \\
  \bottomrule
\end{tabularx}}

\ReportSubsection{Risk-reduction and strategic benefits}
\begin{itemize}
  \item \textbf{CUI containment:} processing stays inside the boundary; no third-party
    retention or training on inputs; fine-tuned models and indexes stay under the
    organization's control.
  \item \textbf{Less shadow AI:} a capable approved tool reduces the incentive to use
    unapproved ones.
  \item \textbf{Auditability:} per-user logs and retrieval records support compliance
    evidence.
  \item \textbf{Continuity:} not dependent on vendor changes to pricing, terms or
    features; a failover path exists for the ML service.
  \item \textbf{Strategic:} builds in-house AI engineering capability, makes historical
    documents searchable, and provides a base for domain-specific assistants.
\end{itemize}

% ===========================================================================
\ReportSection[sec:costs]{Costs}
\TodoAnchor{bj-costs}

Recommended configuration D3; replace all figures with quotes.

\TableTitle{Cost elements, recommended configuration D3, Dell-configured pricing (expected case unless a range is shown)}
{\footnotesize
\begin{xltabular}{\textwidth}{@{}Y{1.4} Y{0.8} Y{0.8}@{}}
  \toprule
  Cost element & ML node(s) (Computer~2) & Computer~1 \\
  \midrule
  \endhead
  Hardware (one-time; street price $\times$ 1.3) & \$64K (street \$49K) & \$20K (street \$15K; \$12--18K) \\
  Setup labor (one-time): ML 200--400~hr, Computer~1 80--160~hr, at \$48.08 & \$14.4K & \$5.8K \\
  \textbf{Upfront cost} & \textbf{\$78K} & \textbf{\$25K} \\
  \midrule
  Non-labor operating cost per year (support, software, power and cooling) & \$7.5K & \$1.8K \\
  Upkeep labor per year: ML 3~hr/week, Computer~1 1~hr/week & \$7.5K & \$2.5K \\
  \textbf{Run cost per year} & \textbf{\$15K} & \textbf{\$4.3K} \\
  \midrule
  \textbf{3-year cost} (upfront $+$ 3 years of run cost) & \textbf{\$123K} (range \$98--145K) & \textbf{\$38K} \\
  \textbf{3-year cost of the whole system (D3)} & \multicolumn{2}{r}{\textbf{\$162K} (range \$131--192K)} \\
  Memo: 3-year cost at street prices & \$108K & \$34K \\
  \bottomrule
\end{xltabular}}
\TodoAnchor{bj-statuscost}
Expansion: about \$8--10K per added GPU; about \$18--50K for an additional node.
If a second Computer~1 is later required for redundancy, add about \$12--18K plus
setup and upkeep. Costs of the status quo to include (fill in): current license and
subscription fees \$\PH{ }/yr, support effort \$\PH{ }/yr, productivity lost to the
limitations, and the cost of any CUI incident or rework.

\begin{ReportNote}{Pricing basis}
Hardware applies a multiplier of 1.3 (range 1.15--1.4) to street prices to approximate
Dell-configured servers; replace with a Dell quote. Windows Server, database and
client-access licensing for Autodesk and Bluebeam server components (about \$2--8K,
estimate) is not included. Labor uses salary only; adding a fringe or overhead
multiplier of 1.3 raises labor figures by about 30\%.
\end{ReportNote}

% ===========================================================================
\ReportSection[sec:fin]{Financial Analysis}

\ReportSubsection{Method}
\[
  \text{Annual benefit} = N_{\text{active}} \times H \times W \times R \times E
\]
where $N_{\text{active}}$ is the number of regular active users, $H$ the hours saved
per user per week, $W$ the working weeks per year (about 46), $R$ the labor rate
(\$48.08/hr) and $E$ a realization factor (the fraction of saved time turned into
useful work; use 0.5--0.7).
\[
  \text{ROI} = \frac{\text{3-yr benefit} - \text{3-yr cost}}{\text{3-yr cost}}
  \qquad
  \text{Payback (months)} = \frac{\text{Upfront cost}}{\text{Net annual benefit}/12}
\]
Cost for the scenarios below is the ML side, which is the increment that depends on
the LLM decision; the last column also attributes Computer~1 to the project.

\ReportSubsection{Illustrative scenarios (all inputs are placeholders)}
\TodoAnchor{bj-scen}

\TableTitle[tbl:scen]{Scenarios}
{\footnotesize
\begin{xltabular}{\textwidth}{@{}Y{1.2} Y{0.95} Y{0.95} Y{0.95} Y{0.95}@{}}
  \toprule
  & Conservative (D3) & Expected (D3) & Optimistic (D4) & Expected, with Computer~1 \\
  \midrule
  \endhead
  Active users $N$ & 20 & 30 & 50 & 30 \\
  Hours saved per user per week $H$ & 1.0 & 2.0 & 3.0 & 2.0 \\
  Working weeks $W$ & 46 & 46 & 46 & 46 \\
  Labor rate $R$ & \$48.08 & \$48.08 & \$48.08 & \$48.08 \\
  Realization $E$ & 0.5 & 0.6 & 0.7 & 0.6 \\
  \textbf{Annual benefit} & \textbf{\$22K} & \textbf{\$80K} & \textbf{\$232K} & \textbf{\$80K} \\
  Annual run cost & \$17K & \$15K & \$16.5K & \$19.3K \\
  Upfront cost (Dell-configured) & \$99K & \$78K & \$91K & \$104K \\
  \textbf{Net annual benefit} & \textbf{\$5K} & \textbf{\$65K} & \textbf{\$216K} & \textbf{\$60K} \\
  \textbf{Payback} & not within 3 years (about 19 years) & \textbf{about 15 months} & about 5 months & about 21 months \\
  \textbf{3-year ROI} & about $-$56\% & \textbf{about 94\%} & about 395\% & about 48\% \\
  \bottomrule
\end{xltabular}}

\begin{ReportNote}{Reading the scenarios}
With only the minimum 20 users each saving 1~hour a week, time savings alone do not
repay the ML cost (conservative column, which also uses the high end of hardware and
setup costs); the case then rests on the capability and risk arguments
(Section~\ref{sec:benefits}). The optimistic case (50 named users, about 30 at once) needs
D4 with both nodes serving and is unlikely as a planning basis. Productivity savings are
often not budget-recoverable ``cash''; present them as \emph{capacity gained} (more
throughput without added headcount) unless a headcount or contract-cost effect can be
shown. Answering from archived standards and past projects may change how much time each
user saves; this is not measured here, and the pilot should measure it.
\end{ReportNote}

\ReportSubsection{Sensitivity and break-even}
\begin{itemize}
  \item \textbf{Biggest driver:} hours saved per user per week. In the expected case,
    each additional 0.5~hr/week across 30 users is worth about \textbf{\$20K per year}.
  \item \textbf{Break-even on run cost:} the ML run cost (\$15K/yr) is recovered with
    about \textbf{11 hours per week} of gross time saved across all users (about 22
    minutes per user for 30 users).
  \item \textbf{Break-even on 3-year cost:} recovering the full 3-year ML cost (\$123K)
    needs about \textbf{31 hours per week} (about 1~hour per user for 30 users); with
    Computer~1 attributed (\$162K), about \textbf{41 hours per week} (at $E = 0.6$).
  \item \textbf{Adoption:} benefits scale with \emph{active} users; plan training and
    champions.
  \item \textbf{Upkeep:} doubling upkeep adds about \$7.5K/yr for D3; upkeep will rise
    during growth phases (adding GPUs, model changes, authorization re-reviews).
  \item \textbf{Utilization:} monitor GPU load; add a GPU at 70\% sustained utilization.
\end{itemize}

\ReportSubsection[sec:node-b]{Is the second ML node worth it?}
\TodoAnchor{bj-nodeb}

Node~B (the difference between S2 and D3) adds about \$20K of hardware at street prices
(about \$26K Dell-configured), 40--80 hours of setup, about 1~hr/week of upkeep and
\$2--3.5K/yr of operating cost: roughly \textbf{\$41--48K over 3 years}.

In the expected case a full day of LLM downtime loses about \$350 of realized benefit
(30 users $\times$ 2~hr/week $\div$ 5 days $\times$ \$48.08 $\times$ 0.6). Failover alone
would therefore take well over 100 avoided outage days to pay for itself. The case for
Node~B rests on its stated main use, training alongside serving, which the productivity
model does not value: archive ingestion, OCR, embedding and fine-tuning can run during
the day without taking GPUs from users. On S2 (2 GPUs, about \$57K Dell-configured),
training during working hours would take a GPU away and leave serving at about 3--8
users, below the stated 12, so training would have to be scheduled off-hours. Node~B can
also act as a second serving node when idle. \PH{Confirm the training workload and the
required recovery time; see Section~\ref{sec:open}.}

\ReportSubsection{Comparison of paths (3-year view)}

\TableTitle{Options compared}
{\footnotesize
\begin{xltabular}{\textwidth}{@{}Y{1.25} Y{0.909} Y{0.795} Y{0.682} Y{1.364}@{}}
  \toprule
  Option & 3-yr cost (est.) & Meets context and function need & CUI control & Comment \\
  \midrule
  \endhead
  Continue status quo & \$\PH{license} & No & Existing & Capability gap persists; shadow-AI risk; no storage or license consolidation. \\
  Upgrade or expand current tools & \$\PH{quote} & Unknown, unlikely & Existing & Ask vendors about context, API and retrieval roadmap. \\
  Compliant cloud LLM $+$ Computer~1 & \$0.07--0.44M (cloud \$0.03--0.4M, usage-driven, $+$ Computer~1 about \$38K) & Yes & Moderate (contract and boundary dependent) & Cost scales with long-context use; data leaves the site; training on the archive would also leave the site. \\
  Computer~1 + S1 (one 1-GPU node) & \$81--121K & Below the stated 12 simultaneous users & Highest & Proof of concept; no failover. \\
  Computer~1 + S2 (one 2-GPU node) & \$94--140K & Yes (training only off-hours) & Highest & No failover. \\
  \textbf{Computer~1 + D3 (recommended)} & \textbf{\$131--192K (ML only about \$98--145K)} & \textbf{Yes} & \textbf{Highest} & Training node that can also serve; reduced-capacity failover. \\
  Computer~1 + D4 & \$144--211K & Yes (more users) & Highest & Full-capacity failover; both nodes can serve. \\
  \bottomrule
\end{xltabular}}
Break-even against a cloud LLM depends on usage: with D3 ML upfront of about \$78K (Dell-configured) and an ML run cost of about \$1.25K/month, cloud spend of \$4K/month breaks even in about
28 months, and the cloud is cheaper if spend stays below about \$3.4K/month (ENG-RPT-01,
Section~\ref{sec:comparison}). For 20--30 users, cloud usage cost alone may be lower; the
on-premises case rests on keeping CUI and the archive inside the boundary, training on
the archive, no per-token cost growth, and the need for Computer~1 in any case.

% ===========================================================================
\ReportSection[sec:usecases]{Use Cases and Target Workflows}
\TodoAnchor{bj-usecases}

Build out with stakeholders.

\TableTitle{Workflows}
{\footnotesize
\begin{xltabular}{\textwidth}{@{}c Y{1.405} c Y{0.649} Y{0.973} Y{0.973}@{}}
  \toprule
  \# & Workflow & Users & Context need & Current tool outcome & Proposed outcome \\
  \midrule
  \endhead
  1 & Specification and requirements review, gap analysis & \PH{ } & 100--300K tokens & Must split; incomplete & Single pass with citations \\
  2 & Test report and data summarization & \PH{ } & 50--150K & Truncated & Complete summary, extracted tables \\
  3 & Archive and standards knowledge search (e.g., NFPA~780) with citations & \PH{ } & Retrieval + 8--32K & Not available & Cited answers across archive and standards \\
  4 & Script and code development and review & \PH{ } & 30--200K & Limited to snippets & Larger-scope analysis \\
  5 & Report and procedure drafting to templates & \PH{ } & 16--64K & Generic output & Template-aware, consistent \\
  6 & Batch classification and extraction & \PH{ } & Per document & Manual & Automated overnight \\
  7 & Fine-tuned assistant for \PH{domain} & \PH{ } & Varies & Not feasible & Adapter trained on Node~B \\
  8 & Proposal and contract support & \PH{ } & 100K+ & Not feasible & Compliance matrix drafting \\
  9 & Lessons-learned mining & \PH{ } & Retrieval & Not feasible & Trend identification \\
  10 & Autodesk, Bluebeam and \PH{other software} hosting; file share and license services & \PH{ } & -- & \PH{current arrangement} & Consolidated on Computer~1 \\
  \bottomrule
\end{xltabular}}

% ===========================================================================
\ReportSection[sec:multi]{Multi-User Operation (Business View)}

\begin{itemize}
  \item One shared platform serves the organization; each user sees only the documents
    and conversations they are authorized for.
  \item \textbf{Fair-use quotas} and priority classes keep interactive users responsive
    while batch jobs and training run off the serving node.
  \item \textbf{Usage dashboards} show adoption and benefit by team (supports showback
    or chargeback).
  \item \textbf{Incremental scaling:} capacity grows by adding GPUs (about \$9K each) or a
    node, without changing the user experience.
  \item \textbf{Limits to communicate:} the stated need is about 12 people generating at the
    same moment. D3's Node~A is rated for 8--20, so 12 fits with margin; others queue
    briefly. Very long (128K-token) requests consume more capacity than short chats. On
    failover to a 1-GPU Node~B, capacity falls to about 3--8 simultaneous users (quotas
    ration it); a 2-GPU Node~B (D4) restores full capacity.
\end{itemize}

% ===========================================================================
\ReportSection[sec:risks]{Risks, Constraints and Mitigations}

\TableTitle{Risks and mitigations}
{\footnotesize
\begin{xltabular}{\textwidth}{@{}Y{0.769}Y{1.231}@{}}
  \toprule
  Risk or concern & Mitigation \\
  \midrule
  \endhead
  Output errors (hallucination) & Citations, human review of anything used in deliverables, policy and training. \\
  Authorization (ATO) timeline, now covering file share and private cloud as well & Early engagement with the security office; one boundary and System Security Plan for both computers. \\
  Operations staffing & Named owner with about 3--4~hr/week protected for upkeep (more during setup and expansion); documented runbooks and a backup person. \\
  Computer 1 is a single point of failure for storage, licenses and the application tier & RAID-Z2 with a hot spare, redundant power supplies, UPS, tested restore; second gateway on Node~B. No second server (sponsor decision); an offline backup copy (about \$2--4K) is recommended because RAID and snapshots share one machine. \\
  A fine-tuned model states wrong clause numbers or requirements & Retrieval with citations is the primary mechanism; evaluate adapters against known answers; human review of anything used in deliverables. \\
  Standards such as NFPA~780 are copyrighted and the license may not allow AI indexing or training & Check license terms before ingesting. \\
  Autodesk or Bluebeam needs internet sign-in or cloud services outside the boundary & Verify the licensing model and self-hosted options with the vendors before the design is frozen. \\
  License vendor restricts virtualization or hardware binding & Confirm terms before the design is frozen; use a dedicated host if required. \\
  Quality gap versus top commercial models & Pilot with side-by-side evaluation on real tasks; use retrieval and tuned prompts. \\
  TAA/OEM sourcing raises hardware cost 15--40\% & Confirm early; obtain OEM quotes. \\
  Single-node configurations have no ML failover & Cold spares, data backups, cloud contingency plan, or choose a dual-node configuration. \\
  Adoption lower than planned & Champions, training, workflow templates, sponsorship. \\
  Model license or provenance & Legal review; approved U.S. or allied-origin models. \\
  Fine-tuned models inherit CUI sensitivity & Registry access control; include in backup and sanitization rules. \\
  Over-reliance on AI & Governance policy and acceptable-use rules. \\
  \bottomrule
\end{xltabular}}

% ===========================================================================
\ReportSection[sec:kpi]{Success Metrics}

\TableTitle{Key performance indicators}
{\small
\begin{xltabular}{\textwidth}{@{}Y{0.724}Y{1.448}Y{0.828}@{}}
  \toprule
  Category & Metric & Target (initial) \\
  \midrule
  \endhead
  Adoption & Weekly active users / named users & 60\% or more by month 6 \\
  Productivity & Hours saved per user per week (survey + measured) & 1.5~hr or more \\
  Quality & User satisfaction & 4.0 / 5 or more \\
  Performance & 95th-percentile time to first token (prompts up to 16K) & Under 3~s \\
  Performance & Streaming speed per user & 15--20 tokens/s or more \\
  Capacity & Sustained GPU utilization on the serving node & 40--70\% \\
  Reliability & Availability of LLM service during duty hours & 98\% or more \\
  Resilience & Failover test passed (dual-node) and restore test passed & Before production; then quarterly \\
  Training & Fine-tuning jobs completed without affecting serving latency & All \\
  Compliance & Open high security findings / CUI incidents & 0 / 0 \\
  Capability & Workflows enabled that were impossible before & 5 or more within 6 months \\
  \bottomrule
\end{xltabular}}

% ===========================================================================
\ReportSection[sec:plan]{Implementation Plan (Summary)}

\begin{enumerate}
  \item \textbf{Months 0--2:} requirements, user survey, baseline usage of Tools A and B,
    storage and license inventory, sourcing rules, quotes, security engagement.
  \item \textbf{Months 1--3:} pilot on one ML node (borrowed or purchased); benchmark
    candidate models, an adapter fine-tune, retrieval over a licensed standard such as
    NFPA~780, and 12 simultaneous users; side-by-side comparison with Tools A and B.
  \item \textbf{Months 2--6:} procurement, authorization documentation, facility
    preparation.
  \item \textbf{Months 5--8:} deploy Computer~1 and Node~A; migrate the file share and
    license servers; integrate sign-on, logging and retrieval; 20--40 pilot users.
  \item \textbf{Months 8--12:} add Node~B (and GPUs as triggers are met); test failover;
    start fine-tuning workloads; measure KPIs.
\end{enumerate}
\emph{Dependent on lead time and authorization pace; not a commitment.}

% ===========================================================================
\ReportSection[sec:baseline]{Baseline Data to Collect Now}
\TodoAnchor{bj-baseline}

\begin{enumerate}
  \item Number of prospective users by team; roles and labor rates.
  \item Usage statistics and satisfaction for Tool A and Tool B.
  \item Documented cases where context limits or missing functionality blocked a task
    (document sizes, time lost).
  \item A time study of 5--10 representative tasks done manually versus with an LLM.
  \item License and support cost of Tool A and Tool B and contract end dates (possible offset).
  \item Known or suspected use of unapproved AI tools (anonymous survey).
  \item Volume of document repositories to index and of data to store (count, size, formats,
    growth).
  \item Inventory of software needing license servers or a private cloud, with vendor
    virtualization rules.
  \item The expected fine-tuning workload: data sets, how often, and how long jobs may take.
  \item Who will operate both computers and their available time.
\end{enumerate}

% ===========================================================================
\ReportSection[sec:rec]{Recommendation}

Proceed with a benchmark and pilot and with procurement planning for a two-computer,
on-premises system authorized for CUI: Computer~1 (about \$20K Dell-configured) for
storage, licensing and the application tier, and ML nodes in configuration \textbf{D3}
(about \$83K total hardware Dell-configured; about \$64K at street prices), built on
4-GPU-capable chassis so either node can grow. Fund Computer~1 and Node~A first and add
Node~B after the pilot confirms demand and the training workload; plan Node~B with two
GPUs (D4) if failover at the full 12 users is required. The investment is justified by
(1) workflows the current tools cannot perform because of context and functionality
limits, (2) controlled handling of CUI, including the NAS archive and fine-tuned models,
(3) consolidated storage and licensing, (4) no per-token cost growth, and (5) a reusable
internal AI platform. At 20--30 users the time-savings return is modest and a compliant
cloud LLM may cost less on usage alone, so funding should be conditional on pilot results
that confirm demand and measured time savings, on a named operations owner, and on a
Dell quote.

% ===========================================================================
\ReportSection[sec:open]{Open Questions for the Sponsor}
\TodoAnchor{bj-open}

Answers received for v0.4: at least 20 named users with about 12 at once; the model must
use and train on NAS archive data (for example NFPA~780); training alongside serving is the
main use of Node~B, with a second-serving-node option; Autodesk and Bluebeam are to be
hosted; storage was estimated at 40--70~TB within 3 years; the labor assumptions are
accepted for now; redundancy for Computer~1 is disks inside the server; Dell is the
vendor; a cloud fallback is desirable but undecided.

\TableTitle{Questions that remain}
{\footnotesize
\begin{xltabular}{\textwidth}{@{}c Y{1.231} Y{0.769}@{}}
  \toprule
  \# & Question & Why it matters \\
  \midrule
  \endhead
  1 & Archive: how large is it now (TB), what formats (scanned PDFs, Office, DWG/Revit, Bluebeam markups), and how are folders and permissions organized? Which documents go first? & Sets ingestion effort and storage; CAD drawings are not text. \\
  2 & Do the licenses for standards such as NFPA~780 allow internal indexing and AI training? & Publisher terms may forbid it. \\
  3 & Autodesk and Bluebeam: which products, named-user or network licensing, does the enclave allow vendor sign-in, and are Vault or Bluebeam collaboration services needed on-site? & Decides whether license servers are needed and what Computer~1 must host. \\
  4 & Server standard: Windows Server with Hyper-V and Active Directory, or Linux with KVM? & Affects Autodesk and Bluebeam hosting, file-share permissions and licensing cost (about \$2--8K). \\
  5 & Dell quotes for Computer~1 and each node configuration; which Dell models support the chosen GPUs? & Replaces the 1.3 planning multiplier. \\
  6 & While Node~A is down, is serving 3--8 users acceptable, or must all 12 be served (Node~B with 2 GPUs, D4)? & Decides between D3 and D4 (about \$12K more at Dell-configured prices). \\
  7 & How fast will named users grow beyond 20? Is Node~B expected to serve as a second node routinely? & Determines when to add GPUs or a node. \\
  8 & Is an offline backup copy wanted in addition to disk redundancy? & About \$2--4K; covers server loss and ransomware. \\
  9 & Which cloud fallback option(s) in ENG-RPT-01 should be refined (none, non-CUI only, CUI-authorized government cloud, warm standby)? & Defines the contingency and its approvals. \\
  10 & What do the current tools cost, and when do their contracts end? (Unknown at present.) & Offsets part of the new cost. \\
  11 & Branding and handling: center name and logo, marking text, document numbers (ENG-RPT-01, ENG-RPT-02), linked intake reference, distribution list. & Currently placeholders on the cover and control block. \\
  \bottomrule
\end{xltabular}}

% ===========================================================================
\ReportSection[sec:glossary]{Glossary}

\TableTitle{Terms used in this report}
{\footnotesize
\begin{xltabular}{\textwidth}{@{}Y{0.55} Y{1.45}@{}}
  \toprule
  Term & Meaning \\
  \midrule
  \endhead
  ATO (Authority to Operate) & Formal approval to run a system, based on its security documentation and testing. \\
  Context window and token & A token is a piece of text, roughly three-quarters of a word. The context window is how many tokens the model can consider at once; 128K tokens is about 300 pages. \\
  CUI & Controlled Unclassified Information: unclassified government information that must be safeguarded. \\
  Dell-configured price & The price of a server built to order by Dell, including qualified parts, warranty and support. Here estimated as the street price $\times$ 1.3. \\
  Failover & Switching to a standby system when the primary one fails. \\
  Fine-tuning & Further training of an existing model on selected examples to change how it behaves, for example its vocabulary or report format. \\
  GPU and ML node & A graphics processing unit does the model's computing. An ML node is a server holding one or more GPUs. \\
  LLM & Large language model: software that reads and writes text, answers questions and drafts documents. \\
  NAS and private cloud & Network-attached storage is shared file storage. The private cloud is a server that runs several virtual machines for specific software inside the organization. \\
  Named and simultaneous users & People who have accounts, and people getting an answer at the same moment. \\
  Open-weight model & A model whose trained weights are published for local use under a license, so it can run without any outside service. \\
  Payback and ROI & Payback is the time for benefits to repay the upfront cost. ROI (return on investment) is the 3-year benefit minus the 3-year cost, divided by the 3-year cost. \\
  Realization factor & The share of saved time that turns into useful extra work (assumed 0.5--0.7). \\
  Retrieval (RAG) & Finding relevant passages in an index of documents and giving them to the model so its answer is based on them and cites them. \\
  TAA and Section 889 & Trade Agreements Act and NDAA Section~889: sourcing rules limiting country of origin and certain vendors for government-related purchases. \\
  \bottomrule
\end{xltabular}}

% ===========================================================================
\ReportSection[sec:approvals]{Approvals}
\TodoAnchor{bj-approve}

\TableTitle{Approvals}
{\small
\begin{tabularx}{\textwidth}{@{}Y{1.455}Y{0.97}Y{0.97}Y{0.605}@{}}
  \toprule
  Role & Name & Decision & Date \\
  \midrule
  Program / department sponsor & & & \\[10pt]
  IT / CISO / ISSM & & & \\[10pt]
  Finance & & & \\[10pt]
  Legal / contracts & & & \\[10pt]
  \bottomrule
\end{tabularx}}


\end{document}
